% Seitenrand: Innen mind. 25mm
% Außenrand: Korrekturrand 50mm,
%   Randnotizen können in den 50mm-Rand hineinragen

\subsubsection{Papierformat}


\subsubsection{Ränder}

Die Seitenränder sollten gemäß DIN~5008 mind. \unit[2,5]{cm} betragen. Bei
Arbeiten, die auf der Außenseite einen Korrekturrand erfordern, sollte der
Korrekturrand \unit[5]{cm} breit sein.

Die Einstellung des Satzspiegels erfolgt mit dem Paket \texttt{geometry}.

\newpage
\setlayoutscale{0.29}
\printheadingsfalse
\printparametersfalse
%\drawdimensionstrue
\currentpage
\begin{minipage}{0.5\linewidth}
\oddpagelayoutfalse
\pagedesign
\end{minipage}%
\begin{minipage}{0.5\linewidth}
\oddpagelayouttrue
\pagedesign
\end{minipage}


\subsubsection{Duplexdruck}

Sofern nicht explizit anders gefordert, sollten Arbeiten beidseitig gedruckt
werden, um Papier zu sparen. Hierzu ist die Klassenoption \texttt{twoside} zu
setzen. Für kürzere Haus- oder Seminararbeiten ist im Zweifel auch einseitiger
Druck statthaft.

\lstinputlisting[language={[LaTeX]{TeX}},style=block,escapechar=\!,caption={Doppelseitiger Textsatz},label={lst:tpl:twoside}]{code/tpl_twoside}

In umfangereicheren Arbeiten (ab Bachelorarbeit aufwärts), sollten neue
Kapitel stets auf einer rechten Seiten beginnen. Die Klasse \texttt{book} wird
vor Beginn eines neuen Kapitels im Bedarfsfall automatisch eine Leerseite
einfügen. Bei der Klasse \texttt{report} ist hierfür die Option
\texttt{openright} notwendig.

\lstinputlisting[language={[LaTeX]{TeX}},style=block,escapechar=\!,caption={Kapitelbeginn auf rechter Seite},label={lst:tpl:openright}]{code/tpl_openright}




\subsubsection{Bindekorrektur}

Studienabschlussarbeiten sind in aller Regel gebunden einzureichen. Meist
kommt eine Kamm- oder Klebebindung zur Anwendung. Da ein Teil des Innenrandes
von der Bindung und dem gebogenen Papier verdeckt wird, muss der Satzspiegel
ein klein wenig enger gestaltet werden. Somit sind die Proportionen der
aufgeschlagenen Arbeit wieder optisch stimmig. Dieser als Bindekorrektur
bezeichnete Abstand~$K$ bestimmt sich dabei aus der Dicke des Buchblocks~$d$
und der Breite der Bindung~$b$.
%
\begin{equation}
K = \dfrac{1}{2} \cdot \max(d, b)
\end{equation}
%
Die Buchblockdicke lässt sich anhand der Grammatur~$G$ und der Anzahl der
Blätter~$n$ berechnen. Dabei ist die Grammatur in der Einheit $\unit{g / m^2}$
in die Größengleichung einzusetzen. Das Ergebnis der Formel entspricht der
Dicke in Millimetern.
%
\begin{equation}
d = \dfrac{n \cdot G}{1000}
\end{equation}
%
Ein Buchblock aus 100~Seiten (das sind 50~Blätter) hat demnach bei Papier der
Stärke $\unit[80]{g / m^2}$ eine Dicke von \unit[4]{mm}. Eine übliche
Kammbindung hat jedoch eine Breite von \unit[6]{mm}, sodass die Bindekorrektur
auf diesen Wert zu vergrößern ist.

Die Bindekorrektur lässt sich über das Paket \texttt{geometry} einstellen.
Listing~\ref{lst:tpl:bcor} zeigt die entsprechende Konfigurationsoption.

\lstinputlisting[language={[LaTeX]{TeX}},style=block,escapechar=\!,caption={Einstellung der Bindekorrektur},label={lst:tpl:bcor}]{code/tpl_bcor}

Bei Ringbindungen muss keine Bindekorrektur erfolgen.

% DTK 2004/1
